\section{Discussion}
\label{sec:discussion}

\paragraph{What process rewards buy.}
In this testbed the main benefit of the monitorability-aware reward is not that it unlocks behaviour outcome reward cannot reach. With a suitable observation, outcome-only agents also find checking policies. The benefit is that it makes those policies \emph{stable} under continued training and \emph{selective}: checks are concentrated on cued steps. The cue bonus does most of this work (\S\ref{sec:ablations}). The cue bonus is also the most domain-specific term, since it rewards the agent for acting on a signal the reward designer already trusts. The general lesson for process supervision is that a process reward which pays for checking where a trusted signal says checking is valuable can make checking reliable. Whether the agent's checking then remains useful depends on whether that signal remains trustworthy.

\paragraph{What process rewards do not buy.}
Process rewards did not fix length generalization. Under absolute position features, the best reward still fell below the never-check baseline on long chains. The bottleneck was what the policy could condition on, not what it was paid for. The same pattern appears in the budget shift, where the slack feature of the invariant encoder goes out of range. For process-supervised reasoning models, this suggests that ``learned to verify on training problems'' and ``will verify appropriately on longer problems'' are separate claims and should be evaluated separately. Positional shortcuts are a plausible way for the second claim to fail even when the first one holds.

\paragraph{Evaluation lessons.}
Three problems in the earlier version of this study each pushed its conclusion in the wrong direction, and none of them was visible in its headline numbers:
\begin{itemize}
  \item a representation confound;
  \item last-checkpoint evaluation of churning policies;
  \item a metric that credited unchecked faults.
\end{itemize}
Small exact testbeds make such problems cheap to find. We recommend reporting learning curves and per-seed variability for RL on reasoning-like tasks, selecting checkpoints only on in-distribution validation, and testing whether out-of-distribution failures move when the representation changes but the reward does not.

\paragraph{Limitations.}
The environment is deliberately minimal. It has at most one fault, checks are exact, the cue is given rather than learned, and the agent is a small memoryless MLP rather than a language model. Process rewards use privileged step correctness, so they represent an idealised process reward model. We study one learner (Double DQN) in one hyperparameter regime plus a small sensitivity sweep. Policy-gradient methods or recurrent agents might churn less, or might learn different allocations. Finally, our invariant encoder hides information that the optimal policy needs on long chains (\S\ref{sec:residual}). The remaining gap is a representation problem we expose but only partly solve.

\section{Conclusion}
In a small budgeted-verification testbed, process rewards made checking behaviour learnable and stable, and made it target suspicious steps. Whether that behaviour carried over to longer chains was decided mostly by the agent's observation encoding, not by its reward. Correcting a positional confound turned an apparent negative result about monitorability-aware RL into a positive but qualified one. We hope the testbed is useful for probing how process supervision shapes \emph{when} learned reasoners check their work.

\paragraph{Reproducibility.}
All code, configs, trained checkpoints, learning curves and per-seed evaluation results are in the accompanying repository (\url{https://github.com/raj200501/UltraRL}). \texttt{make experiments} reruns every experiment in about four hours on a 4-core CPU; \texttt{make paper} rebuilds every figure, table and number in this paper from the saved results.
